\documentclass[12pt,oneside]{book}

% ============================================================
% METADATOS
% ============================================================
\newcommand{\universidad}{Universidad de Buenos Aires}
\newcommand{\facultad}{Facultad de Derecho}
\newcommand{\materia}{Derecho de la Integraci\'on}
\newcommand{\titulo}{Autonom\'ia de la voluntad en los contratos internacionales: teor\'ias, autonom\'ia material y conflictual, y sus l\'imites}
\newcommand{\autor}{Isaac Edgar Camacho Ocampo}
\newcommand{\anio}{2026}

% ============================================================
% PAQUETES
% ============================================================
\usepackage[utf8]{inputenc}
\usepackage[T1]{fontenc}
\usepackage[spanish,provide=*]{babel}
\usepackage[
    top=2.5cm,
    bottom=2.5cm,
    left=3cm,
    right=2.5cm,
    headheight=15pt
]{geometry}
\usepackage{amsmath}
\usepackage{amssymb}
\usepackage{mathtools}
\usepackage{graphicx}
\usepackage{float}
\usepackage{caption}
\usepackage{subcaption}
\usepackage{booktabs}
\usepackage{array}
\usepackage{longtable}
\usepackage{tabularx}
\usepackage{enumitem}
\usepackage{xcolor}
\usepackage[most]{tcolorbox}
\usepackage{fancyhdr}
\usepackage{ragged2e}
\usepackage[
    colorlinks=true,
    linkcolor=blue!50!black,
    citecolor=green!40!black,
    urlcolor=blue!60!black,
    pdftitle={\titulo},
    pdfauthor={\autor},
    pdfsubject={Apuntes universitarios},
    bookmarks=true
]{hyperref}

% ============================================================
% CONFIGURACIÓN
% ============================================================
\graphicspath{
    {./img/}
    {./graficos/}
    {./figuras/}
}
\setcounter{tocdepth}{2}
\setlength{\parindent}{0pt}
\setlength{\parskip}{0.5em}
\setlist[itemize]{
    topsep=0.3em,
    itemsep=0.2em
}
\setlist[enumerate]{
    topsep=0.3em,
    itemsep=0.2em
}
\clubpenalty=10000
\widowpenalty=10000
\captionsetup{font=small,labelfont=bf}

% ============================================================
% COMANDOS
% ============================================================
\newcommand{\abs}[1]{\left|#1\right|}
\newcommand{\norm}[1]{\left\lVert#1\right\rVert}
\newcommand{\vect}[1]{\mathbf{#1}}
\newcommand{\deriv}[2]{\frac{d#1}{d#2}}
\newcommand{\pderiv}[2]{\frac{\partial#1}{\partial#2}}

% ============================================================
% ENTORNOS / CAJAS
% ============================================================
\newtcolorbox{definition}[1]{
    enhanced,
    breakable,
    title={Definici\'on: #1},
    fonttitle=\bfseries,
    colback=blue!3,
    colframe=blue!50!black,
    boxrule=0.8pt,
    arc=2mm
}

\newtcolorbox{important}{
    enhanced,
    breakable,
    title={Importante},
    fonttitle=\bfseries,
    colback=yellow!8,
    colframe=orange!70!black,
    boxrule=0.8pt,
    arc=2mm
}

\newtcolorbox{example}[1]{
    enhanced,
    breakable,
    title={Ejemplo: #1},
    fonttitle=\bfseries,
    colback=green!4,
    colframe=green!40!black,
    boxrule=0.8pt,
    arc=2mm
}

\newtcolorbox{formula}[1]{
    enhanced,
    breakable,
    title={F\'ormula / Regla: #1},
    fonttitle=\bfseries,
    colback=purple!4,
    colframe=purple!50!black,
    boxrule=0.8pt,
    arc=2mm
}

\newtcolorbox{exercise}[1]{
    enhanced,
    breakable,
    title={Ejercicio: #1},
    fonttitle=\bfseries,
    colback=gray!5,
    colframe=gray!60!black,
    boxrule=0.8pt,
    arc=2mm
}

\newtcolorbox{note}{
    enhanced,
    breakable,
    title={Nota},
    fonttitle=\bfseries,
    colback=white,
    colframe=gray!50,
    boxrule=0.6pt,
    arc=2mm
}

% ============================================================
% CONFIGURACIÓN FINAL
% ============================================================
\pagestyle{fancy}
\fancyhf{}
\fancyhead[L]{\nouppercase{\leftmark}}
\fancyhead[R]{\materia}
\fancyfoot[C]{\thepage}
\renewcommand{\headrulewidth}{0.4pt}
\renewcommand{\footrulewidth}{0pt}

\fancypagestyle{plain}{
    \fancyhf{}
    \fancyfoot[C]{\thepage}
    \renewcommand{\headrulewidth}{0pt}
}

\begin{document}

% ============================================================
% PORTADA
% ============================================================
\begin{titlepage}

\thispagestyle{empty}

\begin{center}

\vspace*{1cm}

{\large \textsc{\universidad}}

\vspace{0.2cm}

{\large \textsc{\facultad}}

\vspace{3cm}

{\Huge \textbf{\materia}}

\vspace{1cm}

{\LARGE Autonom\'ia de la voluntad en los contratos internacionales}

\vspace{0.4cm}

{\large Teor\'ias, autonom\'ia material y conflictual, y sus l\'imites}

\vspace{0.8cm}

\textit{Estudiar con constancia es la mejor manera de convertir los conceptos en conocimiento.}

\vspace{1.2cm}

\rule{0.70\textwidth}{0.5pt}

\vspace{0.5cm}

{\large Apuntes de estudio}

\vspace{0.5cm}

\rule{0.70\textwidth}{0.5pt}

\vfill

{\large \autor}

\vspace{0.3cm}

{\large \anio}

\end{center}

\vfill

\begin{flushleft}
\textit{Este es un modesto aporte a los estudiantes de ``Derecho de la Integraci\'on'' no pretende sustituir el material oficial de la materia.}
\end{flushleft}

\end{titlepage}

% ============================================================
% ÍNDICE
% ============================================================
\tableofcontents
\clearpage

% ============================================================
% CONTENIDO
% ============================================================

% ------------------------------------------------------------
\chapter[Autonom\'ia de la voluntad]{Autonom\'ia de la voluntad: concepto y teor\'ias}
% ------------------------------------------------------------

\section{Normas imperativas y normas dispositivas}

% TODO: contenido incompleto o ambiguo en las notas originales (la explicaci\'on comienza ya en curso; el primer fragmento de la clase no consta en el apunte)

La autonom\'ia de la voluntad no puede desplazar a las normas \textbf{imperativas} o \textbf{coactivas}. Opera solamente en un \'ambito de cuestiones donde las normas se consideran \textbf{supletorias} de la voluntad de las partes.

\begin{definition}{Normas dispositivas y normas imperativas}
Las normas \textbf{dispositivas} (o supletorias) \emph{disponen, pero no obligan}: las partes pueden apartarse de ellas y hacer otra cosa. Las normas \textbf{imperativas} o \textbf{coactivas}, en cambio, no pueden ser desplazadas por la autonom\'ia de la voluntad.
\end{definition}

Los negocios y los contratos se encuentran normalmente en el \'ambito de las normas dispositivas. Todos los sistemas jur\'idicos tienen normativa sobre contrato de transporte, contrato de seguro, contrato de compraventa internacional de mercader\'ias y contrato de alquiler; sin embargo, otorgan a las partes un \'ambito para modificar esas cuestiones. Las partes pueden dejar de lado determinadas cuestiones establecidas por la ley.

\section{La postura que niega la autonom\'ia de la voluntad}

En la posici\'on totalmente enfrentada a la de quienes sosten\'ian que la autonom\'ia era absoluta ---algunos afirmaban que ``ven\'ia de Dios''--- se encuentra un sector de la doctrina que sostiene que la autonom\'ia de la voluntad, en realidad, \textbf{no existe}.

El razonamiento es el siguiente. Si la autonom\'ia de la voluntad queda determinada por la existencia de normas facultativas o dispositivas, las partes pueden ejercer su libertad de acci\'on y su libertad contractual \'unicamente mientras se mantengan en ese \'ambito. Si la persona es libre porque el legislador le permite hacer determinada cosa, entonces el \'ambito de su autonom\'ia y de su libre determinaci\'on depende de lo que dispone otro, y la determinaci\'on en realidad ``no es tuya'': el otro est\'a permitiendo hacerlo. De all\'i concluyen que la autonom\'ia de la voluntad no existe como tal.

Lo que existe es el principio de \textbf{libertad de contrataci\'on}, y el \'ambito en el que esta puede desarrollarse \textbf{deriva de la ley}.

% TODO: contenido incompleto o ambiguo en las notas originales (la docente alude a ``las tres teor\'ias'', pero la explicaci\'on inicial de la primera teor\'ia qued\'o fuera de la grabaci\'on)

\begin{important}
Se mencionan tres teor\'ias porque, seg\'un la teor\'ia que adopte cada Estado, se examinan los fundamentos y la extensi\'on que este otorga a la autonom\'ia de la voluntad.
\end{important}

\section{Las teor\'ias y su evoluci\'on hist\'orica}

Las distintas posiciones no son solamente tres teor\'ias: coinciden, adem\'as, con determinadas etapas de evoluci\'on. Cuando surgi\'o la \emph{lex mercatoria} exist\'ia autonom\'ia absoluta. Como referencia literaria de esa concepci\'on se menciona \emph{El mercader de Venecia}, de Shakespeare, lo que indica que no se trataba de una idea extra\~na.

Esta coincidencia con un desarrollo hist\'orico responde, ya sea a una convicci\'on jur\'idica, ya sea a una conveniencia de los negocios, que lleva a adoptar o modificar las teor\'ias. Hoy, doctrinariamente, por m\'as que un pa\'is declare adoptar o no la autonom\'ia de la voluntad, \textbf{nadie niega} de modo tajante la noci\'on de autonom\'ia de la voluntad. Se habla de autonom\'ia de la voluntad y no de libertad de contrataci\'on.

\begin{table}[H]
\centering
\caption{Las posiciones sobre la autonom\'ia de la voluntad (ordenado a partir de lo dicho en clase)}
\small
\begin{tabularx}{\textwidth}{
    >{\raggedright\arraybackslash}p{0.17\textwidth}
    >{\raggedright\arraybackslash}X
    >{\raggedright\arraybackslash}X
    >{\raggedright\arraybackslash}X
}
\toprule
\textbf{Posici\'on} & \textbf{Qu\'e sostiene} & \textbf{Fundamento} & \textbf{Etapa / observaci\'on} \\
\midrule
\textbf{Autonom\'ia absoluta} &
Las partes pueden pactar sin l\'imites; lo pactado rige tal cual. &
Algunos decían que ``ven\'ia de Dios''. &
Se asocia a la \emph{lex mercatoria}; ejemplo literario: \emph{El mercader de Venecia}. \\
\addlinespace
\textbf{Negaci\'on de la autonom\'ia} &
No existe autonom\'ia; solo existe libertad de contrataci\'on. &
La libertad se ejerce en el \'ambito que el legislador permite (normas dispositivas). &
Grupo de la doctrina ``totalmente enfrentado'' a la postura absoluta. \\
\addlinespace
\textbf{Autonom\'ia limitada (actual)} &
Existe autonom\'ia de la voluntad, pero con l\'imites. &
Nadie niega hoy la noci\'on; se habla de autonom\'ia de la voluntad y no de libertad de contrataci\'on. &
Posici\'on actual: l\'imites en las estipulaciones y en la ley que se puede elegir. \\
\bottomrule
\end{tabularx}
\end{table}

\begin{note}
El cuadro ordena las posiciones mencionadas en la clase. La explicaci\'on inicial de la primera teor\'ia qued\'o fuera de la grabaci\'on.
\end{note}

\section{Cuadro Cornell}

\begin{longtable}{
    >{\raggedright\arraybackslash}p{0.26\textwidth}
    >{\raggedright\arraybackslash}p{0.66\textwidth}
}
\toprule
\textbf{Preguntas / palabras clave} & \textbf{Notas / respuestas} \\
\midrule
\endfirsthead
\toprule
\textbf{Preguntas / palabras clave} & \textbf{Notas / respuestas} \\
\midrule
\endhead
\bottomrule
\endfoot
\bottomrule
\endlastfoot

\textbf{¿Qu\'e diferencia hay entre norma imperativa y norma dispositiva?} &
Las imperativas (coactivas) no pueden ser desplazadas por la autonom\'ia de la voluntad. Las dispositivas (supletorias) ``disponen, no te obligan'': las partes pueden dejarlas de lado. Los negocios y contratos se mueven normalmente en este segundo \'ambito. \\
\addlinespace

\textbf{¿Qu\'e sostiene la postura que niega la autonom\'ia de la voluntad?} &
Que, si la libertad de las partes depende de lo que el legislador les permite, la determinaci\'on ``no es tuya''. Lo que existe es el principio de libertad de contrataci\'on, cuyo \'ambito deriva de la ley. \\
\addlinespace

\textbf{¿Para qu\'e importa conocer las teor\'ias?} &
Porque, seg\'un la teor\'ia que adopte un Estado, se examinan los fundamentos y la extensi\'on que da a la autonom\'ia de la voluntad. Adem\'as, las teor\'ias coinciden con etapas hist\'oricas (\emph{lex mercatoria}, autonom\'ia absoluta; hoy, autonom\'ia limitada). \\
\addlinespace

\textbf{¿Qu\'e postura rige hoy?} &
Nadie niega la noci\'on de autonom\'ia de la voluntad, aunque limitada. Se habla de autonom\'ia de la voluntad y no de libertad de contrataci\'on. \\

\midrule
\multicolumn{2}{>{\raggedright\arraybackslash}p{\dimexpr0.92\textwidth+2\tabcolsep\relax}}{%
\textbf{Resumen:} las normas dispositivas permiten a las partes apartarse de la ley y las imperativas no; en ese \'ambito supletorio se mueven los contratos. Una posici\'on doctrinaria niega la autonom\'ia y la reduce a libertad de contrataci\'on derivada de la ley; las teor\'ias acompa\~nan etapas hist\'oricas, desde la autonom\'ia absoluta de la \emph{lex mercatoria} hasta la actual autonom\'ia limitada.} \\
\end{longtable}

% ------------------------------------------------------------
\chapter[Autonom\'ia material y conflictual]{Autonom\'ia material y autonom\'ia conflictual}
% ------------------------------------------------------------

\section{Autonom\'ia material}

La autonom\'ia material es la libertad de contrataci\'on: las estipulaciones insertas en el contrato son, para las partes, una regla a la que deben someterse como a la ley misma. Brasil, por ejemplo, acepta la autonom\'ia material. Esa idea corresponde a la formulaci\'on de V\'elez Sarsfield.

\begin{formula}{Art. 1197 del C\'odigo Civil de V\'elez Sarsfield (derogado)}
\emph{<<Las convenciones hechas en los contratos forman para las partes una regla a la cual deben someterse como a la ley misma, y no pueden ser revocadas sino por mutuo consentimiento o por las causas que la ley autoriza.>>}

\smallskip
\small Norma aludida en la clase. Texto del C\'odigo Civil de la Naci\'on vigente hasta 2015.
\end{formula}

\begin{formula}{Art. 959 del C\'odigo Civil y Comercial de la Naci\'on (norma vigente concordante)}
\emph{<<Todo contrato v\'alidamente celebrado es obligatorio para las partes. Su contenido solo puede ser modificado o extinguido por acuerdo de partes o en los supuestos en que la ley lo prev\'e.>>}

\smallskip
\small Agregado del apunte original: no fue mencionado en la clase, pero es la norma que hoy recoge la misma idea.
\end{formula}

Casi todos los ordenamientos jur\'idicos autorizan a las partes a celebrar contratos y establecen que sus obligaciones quedar\'an reguladas en la medida en que lo hayan establecido en ellos.

\begin{example}{Qu\'e d\'ia se entrega, qu\'e precio se paga}
¿Qu\'e d\'ia est\'a obligada a entregar la mercader\'ia una de las partes? El d\'ia fijado en el contrato; eso es obligatorio para ella. ¿Qu\'e precio debe pagar? Lo que se puso en el contrato. Esas estipulaciones son obligatorias. Eso es \textbf{autonom\'ia material}, lo que otros llaman libertad de contrataci\'on.
\end{example}

Para quienes niegan la autonom\'ia, esto no es autonom\'ia, porque la autonom\'ia \emph{per se} no existe. Hoy se sostiene que la autonom\'ia de la voluntad existe, pero limitada. Existir\'an determinados l\'imites, tanto en cuanto a las estipulaciones que se incluyen en el contrato ---qu\'e se puede pactar y qu\'e no--- como en cuanto a la ley que puede elegirse para regular ese contrato. Esto \'ultimo es la autonom\'ia conflictual.

\section{Autonom\'ia conflictual}

Las cl\'ausulas que las partes incluyen en el contrato constituyen la autonom\'ia material. En cambio, si en el contrato se pacta, por ejemplo: <<El presente contrato, los derechos y obligaciones, responsabilidades, la interpretaci\'on de sus cl\'ausulas y dem\'as, ser\'a regido por la ley uruguaya>>, las partes establecieron qu\'e ley dir\'a si lo que hicieron en el contrato est\'a bien.

% TODO: contenido incompleto o ambiguo en las notas originales (el ejemplo del precio excesivo y del inter\'es usurario se enuncia de forma incompleta)

Se pod\'ia celebrar ese contrato de compraventa: el objeto era l\'icito, el precio era excesivo, el inter\'es era usurario. Conforme a la ley que va a regular si esas estipulaciones son v\'alidas, y la interpretaci\'on de ellas, ya no se trata de una autonom\'ia absoluta.

\begin{example}{El mercader de Venecia}
Si se pact\'o la entrega de una parte del coraz\'on, debe cumplirse. La objeci\'on de que ello viola el orden p\'ublico no tuvo lugar: no se habl\'o de orden p\'ublico ni se accedi\'o a esas soluciones. El caso se resolvi\'o conforme a los mismos t\'erminos del contrato.
\end{example}

\begin{table}[H]
\centering
\caption{Autonom\'ia material y autonom\'ia conflictual}
\small
\begin{tabularx}{\textwidth}{
    >{\raggedright\arraybackslash}p{0.19\textwidth}
    >{\raggedright\arraybackslash}X
    >{\raggedright\arraybackslash}X
}
\toprule
 & \textbf{Autonom\'ia material} & \textbf{Autonom\'ia conflictual} \\
\midrule
\textbf{Qu\'e es} & Libertad de fijar el contenido del contrato: las cl\'ausulas. & Posibilidad de elegir la ley que regir\'a el contrato. \\
\addlinespace
\textbf{Otro nombre} & Libertad de contrataci\'on. & --- \\
\addlinespace
\textbf{Ejemplo de la clase} & Fecha de entrega y precio fijados en el contrato. & <<Este contrato ser\'a regido por la ley uruguaya>>. \\
\addlinespace
\textbf{Aceptaci\'on} & Todos la aceptan. & Va <<a otro nivel>>: ning\'un Estado la admite de manera absoluta. \\
\addlinespace
\textbf{L\'imites} & Qu\'e se puede pactar y qu\'e no (p.\,ej., orden p\'ublico). & La ley elegida puede ser desplazada si afecta principios fundamentales del Estado. \\
\bottomrule
\end{tabularx}
\end{table}

\section{Qu\'e ley rige cada aspecto del contrato}

Hay tres cuestiones distintas, cada una regida por su ley. En primer lugar, la \textbf{celebraci\'on del contrato}: oferta, demanda, momento y lugar de celebraci\'on; esto determina si existe contrato. En segundo lugar, otra ley regula la \textbf{validez extr\'inseca}, es decir, las formalidades del acto. En tercer lugar, otra normativa determina si el contrato es \textbf{sustancialmente v\'alido}.

Superadas la celebraci\'on y las formalidades, corresponde examinar el objeto. Aunque el contrato se haya hecho por escribano p\'ublico y tenga oferta y precio, puede ocurrir que el objeto no estuviera a la venta, que no fuera susceptible de comercializaci\'on (se da como ejemplo la venta del Obelisco) o que se tratara de un contrato de compraventa internacional de marihuana.

\begin{table}[H]
\centering
\caption{Las tres cuestiones del contrato y su ley}
\small
\begin{tabularx}{\textwidth}{
    >{\raggedright\arraybackslash}p{0.22\textwidth}
    >{\raggedright\arraybackslash}X
    >{\raggedright\arraybackslash}X
}
\toprule
\textbf{Cuesti\'on} & \textbf{Qu\'e determina} & \textbf{Ejemplo de la clase} \\
\midrule
\textbf{Celebraci\'on del contrato} & Si el contrato existe: oferta, demanda, momento y lugar de celebraci\'on. & Hubo oferta, hubo precio. \\
\addlinespace
\textbf{Validez extr\'inseca (formal)} & Las formalidades del acto. & El contrato se hizo por escribano p\'ublico. \\
\addlinespace
\textbf{Validez sustancial} & Si el contrato es sustancialmente v\'alido: objeto, precio, inter\'es. & El objeto no era susceptible de comercializaci\'on (el Obelisco; la marihuana). \\
\bottomrule
\end{tabularx}
\end{table}

\begin{example}{Compraventa internacional de marihuana, ley uruguaya y juez argentino}
Se celebra un contrato de compraventa con Uruguay, formalmente escrito, y se designa como ley aplicable la ley uruguaya. La ley uruguaya permite la compraventa de marihuana para consumo personal. El juez uruguayo que aplique su ley considerar\'a v\'alido el contrato: hubo autonom\'ia de la voluntad de las partes y, mientras sea para uso personal y por la cantidad adquirida, no hay problema.

Si el juez es argentino, al ir a aplicar la ley uruguaya dir\'a que ello va contra el orden p\'ublico argentino, porque el objeto es il\'icito. No aplicar\'a la ley uruguaya y declarar\'a que esa elecci\'on de ley no es v\'alida. La raz\'on es que la elecci\'on de la ley aplicable hecha por autonom\'ia de la voluntad tiene ciertos l\'imites: no se puede hacer cualquier cosa ni elegir cualquier ley.
\end{example}

M\'as all\'a de las estipulaciones del contrato, es necesario saber qu\'e ley analizar\'a si el contrato es ejecutable o no. Cuando se afirma que una cl\'ausula es nula, se la est\'a juzgando conforme a una determinada ley; de ah\'i la pregunta por la ley aplicable al contrato.

Si se parte de que el Estado adopta la autonom\'ia conflictual, es decir, la posibilidad de que las partes elijan la ley aplicable, ning\'un Estado la adopta de manera absoluta (``elijo la ley que quiero''). Quienes aceptan la autonom\'ia conflictual establecen normalmente algunos l\'imites. El primero es el orden p\'ublico: si la ley elegida va contra el orden p\'ublico del Estado, no se aplica.

\begin{important}
M\'as all\'a de las cl\'ausulas, hace falta saber qu\'e ley va a juzgar si el contrato es v\'alido y ejecutable. Un mismo contrato puede ser v\'alido para el juez uruguayo y no ser reconocido por el juez argentino, porque la elecci\'on de la ley tiene l\'imites: ninguna ley elegida se aplica si va contra el orden p\'ublico del Estado del juez.
\end{important}

\section{Cuadro Cornell}

\begin{longtable}{
    >{\raggedright\arraybackslash}p{0.26\textwidth}
    >{\raggedright\arraybackslash}p{0.66\textwidth}
}
\toprule
\textbf{Preguntas / palabras clave} & \textbf{Notas / respuestas} \\
\midrule
\endfirsthead
\toprule
\textbf{Preguntas / palabras clave} & \textbf{Notas / respuestas} \\
\midrule
\endhead
\bottomrule
\endfoot
\bottomrule
\endlastfoot

\textbf{¿Qu\'e es la autonom\'ia material?} &
La libertad de fijar las cl\'ausulas del contrato (fecha de entrega, precio). Lo estipulado obliga a las partes como la ley misma (art.~1197 del C\'odigo Civil derogado; hoy, art.~959 CCCN, seg\'un el agregado del apunte). Otros la llaman libertad de contrataci\'on. \\
\addlinespace

\textbf{¿Qu\'e es la autonom\'ia conflictual?} &
La posibilidad de que las partes elijan la ley que regir\'a el contrato, su validez y la interpretaci\'on de sus cl\'ausulas (por ejemplo: <<ser\'a regido por la ley uruguaya>>). Ning\'un Estado la admite de forma absoluta. \\
\addlinespace

\textbf{¿Qu\'e leyes intervienen en un contrato?} &
Una regula la celebraci\'on (oferta, demanda, momento y lugar); otra la validez extr\'inseca (formalidades); otra la validez sustancial (objeto, precio, inter\'es). \\
\addlinespace

\textbf{¿Qu\'e pasa con el contrato de compraventa de marihuana con ley uruguaya?} &
El juez uruguayo, aplicando su ley, lo considera v\'alido. El juez argentino no aplica la ley elegida porque va contra el orden p\'ublico argentino: objeto il\'icito. La elecci\'on de ley tiene l\'imites. \\

\midrule
\multicolumn{2}{>{\raggedright\arraybackslash}p{\dimexpr0.92\textwidth+2\tabcolsep\relax}}{%
\textbf{Resumen:} la autonom\'ia material consiste en fijar las cl\'ausulas del contrato, y todos la aceptan; la autonom\'ia conflictual consiste en elegir la ley aplicable y ning\'un Estado la admite de manera absoluta. Un mismo contrato, como la compraventa de marihuana regida por ley uruguaya, puede ser v\'alido para un juez y no ser reconocido por otro, porque la elecci\'on de ley tiene l\'imites, el primero de ellos el orden p\'ublico.} \\
\end{longtable}

% ------------------------------------------------------------
\chapter[L\'imites: orden p\'ublico y normas de polic\'ia]{L\'imites: orden p\'ublico y normas de polic\'ia}
% ------------------------------------------------------------

\section{El orden p\'ublico como l\'imite}

Los l\'imites m\'as habituales que aparecen en los distintos ordenamientos jur\'idicos a la autonom\'ia de la voluntad conflictual ---es decir, a la posibilidad de elegir la ley aplicable--- incluyen el \textbf{orden p\'ublico}.

\begin{definition}{Orden p\'ublico}
Es algo tan importante para el Estado que no admite la autonom\'ia de la voluntad de las partes y puede frenar la aplicaci\'on del derecho extranjero. Si lo pactado, o el resultado de aplicar la ley elegida, afecta principios fundamentales del Estado, no se reconoce.
\end{definition}

Cuando se verifica que el acto conculca esos principios fundamentales, no se aplica ni la autonom\'ia material ni la autonom\'ia conflictual. Se distingue entre orden p\'ublico interno y orden p\'ublico internacional.

Como ejemplo se mencion\'o el matrimonio de tres personas, permitido en algunos pa\'ises: nuestro orden p\'ublico no lo aceptar\'ia. Distinto es el caso de la compraventa de marihuana, que se plantea en materia de negocios: aqu\'i hay un delito, y aunque algunos sostengan que para consumo est\'a autorizado, la cuesti\'on exige analizar lo que se expone a continuaci\'on.

\section{Art\'iculo 19 y consumo personal}

Se sostiene que la Corte autoriz\'o el consumo en el famoso fallo ``para consumo personal''. Esa es una interpretaci\'on muy libre, propia de quienes leyeron superficialmente el fallo: \textbf{la Corte no lo autoriz\'o}. La Corte se\~nal\'o que hay principios, y uno de ellos es la libertad, consagrada en el art\'iculo 19.

\begin{formula}{Art. 19 de la Constituci\'on de la Naci\'on Argentina}
\emph{<<Las acciones privadas de los hombres que de ning\'un modo ofendan al orden y a la moral p\'ublica, ni perjudiquen a un tercero, est\'an solo reservadas a Dios, y exentas de la autoridad de los magistrados. Ning\'un habitante de la Naci\'on ser\'a obligado a hacer lo que no manda la ley, ni privado de lo que ella no proh\'ibe.>>}

\smallskip
\small La docente lo cita de memoria; el apunte original transcribe el texto constitucional.
\end{formula}

\begin{example}{La Corte y la tenencia para consumo personal}
Se trata de situaciones en las que se detiene a determinados j\'ovenes, en circunstancias sospechosas, se les revisan los bolsillos y se les encuentra droga. Es algo interno de la persona, como si se allanara una casa y se encontrara un cigarrillo de marihuana en la mesita de luz. La Corte dijo que, si es para consumo personal, la violaci\'on del derecho a la intimidad y a las acciones personales es demasiado grave para combatir un delito tan menor.

Si en cambio se consume en la v\'ia p\'ublica, la polic\'ia no est\'a invadiendo el \'ambito de intimidad. No se priva de la intimidad a quien lo hace p\'ublicamente, y puede ser detenido. Lo que la Corte dijo es que no era posible, para perseguir el delito de consumo, invadir \'ambitos privados: ni el domicilio ni la intimidad personal; no se puede abrir la mochila. Distinta es la situaci\'on si se sospecha que en esa casa hay venta: en ese caso se puede allanar, y hay que pedir la orden.

La pol\'itica de la ley anterior para combatir el delito era perseguirlo en todos los grados y en todas las instancias. Despu\'es de 10 a\~nos de vigencia de la norma, se advirti\'o que se estaba procesando a consumidores, a personas que a veces est\'an enfermas y son consumidores porque ya tienen una adicci\'on, y que el tr\'afico no hab\'ia disminuido. Se estaba atacando al eslab\'on d\'ebil, sin producir el beneficio que la sociedad necesita. El delito sigue siendo delito; la Corte dijo que se busque combatirlo por otro lado. La premisa utilizada era que no habr\'ia tr\'afico si no hubiera consumo, y esa hip\'otesis no dio resultado.

En conclusi\'on, la Corte no autoriz\'o la invasi\'on de la propiedad o de la intimidad cuando el delito perseguido era el consumo. En la Argentina no est\'a permitido el consumo: puede hacerse si no se es descubierto.
\end{example}

\begin{note}
La docente no menciona el nombre del fallo; el apunte original indica que se trata de la doctrina de la Corte Suprema sobre tenencia para consumo personal y el art.~19 de la Constituci\'on (<<Arriola>>, 2009).
\end{note}

En Uruguay, en cambio, el consumo personal est\'a permitido; tambi\'en en Pa\'ises Bajos. Si las partes eligen esa ley, el juez que la aplique considerar\'a que el contrato es v\'alido, aunque el vendedor haya sido argentino, y lo har\'a ejecutable (por ejemplo, para reclamar el cumplimiento si no se entreg\'o o por la calidad de lo entregado). La Rep\'ublica Argentina no podr\'ia hacerlo bajo ninguna circunstancia, aunque se haya hecho uso de la autonom\'ia de la voluntad y se haya elegido una ley aplicable: no puede aplicar esa ley por el principio del orden p\'ublico, porque la persecuci\'on de un delito limita la autonom\'ia de la voluntad de las partes.

Existe entonces la autonom\'ia material, con la que todos est\'an de acuerdo, y la autonom\'ia conflictual, que va a otro nivel: elegir la ley que juzgar\'a las disposiciones del contrato.

\section{Normas de polic\'ia y normas coactivas}

Existe la autonom\'ia con ciertos l\'imites. La \textbf{norma de polic\'ia} es parecida al orden p\'ublico, pero m\'as r\'igida. El orden p\'ublico afecta principios fundamentales de la legislaci\'on: su esp\'iritu, no su letra. Con el orden p\'ublico se hace un camino de ida y vuelta: se mira el derecho extranjero y se lo confronta con el propio ordenamiento para ver si hay una confrontaci\'on.

\begin{definition}{Norma de polic\'ia}
Es una norma de \textbf{aplicaci\'on inmediata}. A pesar de que exista un elemento de extranjer\'ia que vincule la relaci\'on jur\'idica con otros ordenamientos, existen cuestiones que el Estado considera que nunca pueden quedar sometidas a una ley extranjera ni a la autonom\'ia de la voluntad de las partes: son cuestiones de aplicaci\'on directa e inmediata, aunque haya elementos de extranjer\'ia.
\end{definition}

Las normas de polic\'ia suelen confundirse con las ``normas coactivas''. Las normas coactivas son normas obligatorias, y hay muchas normas obligatorias que no son normas de polic\'ia. Por ejemplo, el homicidio est\'a prohibido, pero no es una norma de polic\'ia: se aplica la ley del lugar de perpetraci\'on del delito, y no siempre la ley argentina.

Las normas de polic\'ia hacen a la esencia del Estado, con tal importancia e impacto que se necesita algo m\'as fuerte que una norma coactiva, obligatoria o imperativa: algo que se aplique de manera directa y no sea cuestionable. Estas situaciones no son habituales; pueden darse en forma extrema. La idea es: <<esto se hace as\'i y punto>>. Un ejemplo es la pandemia.

\begin{table}[H]
\centering
\caption{Norma coactiva, orden p\'ublico y norma de polic\'ia}
\small
\begin{tabularx}{\textwidth}{
    >{\raggedright\arraybackslash}p{0.15\textwidth}
    >{\raggedright\arraybackslash}X
    >{\raggedright\arraybackslash}X
    >{\raggedright\arraybackslash}X
}
\toprule
 & \textbf{Norma coactiva / imperativa} & \textbf{Orden p\'ublico} & \textbf{Norma de polic\'ia} \\
\midrule
\textbf{Idea central} &
Norma obligatoria: no puede ser desplazada por las partes. &
Principios fundamentales del ordenamiento (<<el esp\'iritu>>, no la letra). &
Cuestiones esenciales para el Estado, de aplicaci\'on directa e inmediata. \\
\addlinespace
\textbf{C\'omo opera} &
Se aplica la ley que corresponda (p.\,ej., homicidio: ley del lugar del delito). &
Camino de ida y vuelta: se mira el derecho extranjero y se lo confronta con el propio. &
No se confronta: se aplica directamente, aunque haya elementos de extranjer\'ia. \\
\addlinespace
\textbf{Rigidez} &
Hay much\'isimas. &
Menos r\'igido que la norma de polic\'ia. &
M\'as r\'igida que el orden p\'ublico; situaciones excepcionales y extremas. \\
\bottomrule
\end{tabularx}
\end{table}

\begin{important}
No toda norma coactiva es una norma de polic\'ia: <<hay un mont\'on de normas obligatorias que no son normas de polic\'ia>>. Las normas de polic\'ia se reservan para lo que hace a la esencia del Estado.
\end{important}

\section{Ejemplo de norma de polic\'ia: la pandemia}

% TODO: contenido incompleto o ambiguo en las notas originales (algunos pasajes de esta explicaci\'on tienen audio defectuoso y se reconstruyeron por el sentido)

\begin{example}{La pandemia}
La pandemia vulner\'o muchos derechos: el derecho al trabajo, el derecho a circular. Fue como si cambiaran las reglas del juego: quien ten\'ia un contrato celebrado y deb\'ia viajar no pod\'ia hacerlo. La ley del lugar de cumplimiento del contrato dec\'ia una cosa, pero las soluciones fueron otras: ``no viaj\'as, y devuelven la plata'' o ``no se devuelve nada, porque si no las empresas quiebran''. No se estableci\'o sin criterio: la situaci\'on era tan extraordinaria que necesitaba una soluci\'on espec\'ifica para ese tema.

Cada Estado dict\'o sus normas. Todos hablaban de que no se pod\'ia transmitir ni contagiar a las personas, pero cada uno lo implement\'o de una manera distinta.
\end{example}

\section{Ejemplo de norma de polic\'ia: la crisis de 2001 y el corralito}

\subsection{C\'omo funcionaba la paridad y por qu\'e faltaban los d\'olares}

\begin{example}{Los dep\'ositos en d\'olares que no exist\'ian}
El Estado hab\'ia fijado la paridad un peso, un d\'olar: eran equivalentes. Quien depositaba un peso en el banco depositaba el equivalente a un d\'olar. Quien ten\'ia una cuenta en d\'olares y depositaba 100, figuraba con 100 d\'olares depositados, cosa que hac\'ia todo el mundo porque se hab\'ia equiparado el valor. No era posible depositar pesos y comprar libremente d\'olares; equival\'ia al d\'olar.

Todo el mundo ten\'ia dep\'ositos en d\'olares sin que hubiera d\'olares en los bancos: en las cajas hab\'ia pesos. Si todos iban juntos a retirar sus d\'olares, no exist\'ian.
\end{example}

\subsection{El corralito y la pesificaci\'on de los contratos}

A la crisis econ\'omica, producto del endeudamiento del Estado, se sum\'o la decisi\'on de establecer el \textbf{corralito financiero}: nadie pod\'ia sacar dinero del banco, ni d\'olares ni pesos. Se autorizaban 200 pesos, por persona, sin importar si la persona era soltera o casada con hijos. Era arbitrario.

Respecto de los contratos pactados en d\'olares, se dispuso que se pagar\'ian en pesos, a una tasa establecida por el Estado. Primero se habl\'o de un peso con 25 (la docente no recuerda con certeza la cifra) y luego se fij\'o 1,40. Si en el contrato se hab\'ia pactado pagar la mercader\'ia 200 d\'olares, deb\'ian pagarse 200 por 1,40 en pesos. Era totalmente arbitrario, porque el valor no estaba fijado por el mercado; no pod\'ia dejarse al valor de mercado porque todo el mundo se desesperar\'ia por conseguir d\'olares, dado que la crisis hab\'ia depreciado el peso enormemente.

El Estado prohibi\'o sacar dinero de los bancos e inform\'atiz\'o todo: todos los salarios deb\'ian pagarse por banco. Antes pasaba el cami\'on de caudales y se pagaba a la gente en la puerta de la oficina; desde entonces se hace por transferencia bancaria. Todo eso surgi\'o porque el dinero era ficticio: no hab\'ia ni d\'olares ni pesos en el banco, porque el Estado se hab\'ia endeudado y hab\'ia tomado esos pesos. Entonces se declar\'o el \emph{default}: se dej\'o de pagar porque ya no hab\'ia dinero.

\subsection{Por qu\'e eran normas de polic\'ia}

Todas las normas dictadas durante ese per\'iodo se presentaban como normas que afectan el orden p\'ublico, pero realmente eran \textbf{normas de polic\'ia}, porque no admit\'ian pacto en contrario ni la aplicaci\'on de un derecho extranjero. Si en el contrato se hab\'ia pactado que el lugar de pago ser\'ia Francia, no pod\'ia pagarse en d\'olares en Francia: se pagaba en pesos, porque no hab\'ia d\'olares para transferirlos.

Por eso muchos acudieron a la Corte Internacional y hubo tantos fallos en contra en esa fecha: para ellos la ley de emergencia no importaba; si las partes hab\'ian pactado algo, el contrato deb\'ia cumplirse en los t\'erminos pactados. No importaba que el Estado hubiera considerado ficticiamente que el valor era el mismo: los d\'olares no exist\'ian, y se ten\'ian certificados de dep\'osito que dec\'ian que se hab\'ian depositado tantos d\'olares.

La crisis fue de tal envergadura que las normativas se dictaban con enorme frecuencia. Lleg\'o al punto de que una norma prohibi\'o la quiebra de las empresas: no se pod\'ian iniciar juicios de quiebra y ning\'un juez pod\'ia decretar una quiebra. En la peor crisis econ\'omica del pa\'is, en cinco a\~nos no quebr\'o ninguna empresa. El motivo era que hacerlo habr\'ia agravado el desastre: cierre de empresas, desempleo y otras consecuencias.

Esas disposiciones eran normas de polic\'ia que cambiaban los contratos de trabajo, los contratos de compraventa y el contrato de dep\'osito bancario; cambiaban todas las reglas del juego. Frente a la pretensi\'on de que se aplicara la ley vigente al celebrarse el contrato, o de que los contratos se rigieran por otro derecho, la respuesta era que se trataba de una norma de polic\'ia y se aplicaba s\'i o s\'i.

% TODO: contenido incompleto o ambiguo en las notas originales (este \'ultimo pasaje tiene audio defectuoso; se reconstruy\'o por el sentido)

\begin{important}
Las normas de emergencia de 2001 se presentaban como <<de orden p\'ublico>>, pero la docente las califica como normas de polic\'ia: no admit\'ian pacto en contrario ni aplicaci\'on de derecho extranjero. Cambiaron los contratos de trabajo, de compraventa y de dep\'osito bancario, y se aplicaron aunque las partes hubieran elegido otra ley o pactado otro lugar de pago.
\end{important}

\section{Cuadro Cornell}

\begin{longtable}{
    >{\raggedright\arraybackslash}p{0.26\textwidth}
    >{\raggedright\arraybackslash}p{0.66\textwidth}
}
\toprule
\textbf{Preguntas / palabras clave} & \textbf{Notas / respuestas} \\
\midrule
\endfirsthead
\toprule
\textbf{Preguntas / palabras clave} & \textbf{Notas / respuestas} \\
\midrule
\endhead
\bottomrule
\endfoot
\bottomrule
\endlastfoot

\textbf{¿C\'omo funciona el orden p\'ublico?} &
Frena la aplicaci\'on del derecho extranjero (o del pacto) cuando su resultado afecta principios fundamentales del Estado. Se hace un camino de ida y vuelta: se mira el derecho extranjero y se lo confronta con el propio. Hay orden p\'ublico interno e internacional. \\
\addlinespace

\textbf{¿Qu\'e dijo la Corte sobre el consumo personal y el art.~19?} &
No autoriz\'o el consumo: dijo que no se puede invadir el domicilio ni la intimidad para perseguir el delito de consumo. En la v\'ia p\'ublica el consumo sigue siendo delito. Para investigar venta s\'i se puede allanar, con orden. \\
\addlinespace

\textbf{¿Qu\'e es una norma de polic\'ia?} &
Una norma de aplicaci\'on inmediata y directa, que el Estado considera esencial y que no admite someterse a ley extranjera ni a la autonom\'ia de la voluntad, aunque haya elementos de extranjer\'ia. Es m\'as r\'igida que el orden p\'ublico. \\
\addlinespace

\textbf{¿Toda norma coactiva es una norma de polic\'ia?} &
No. Hay muchas normas obligatorias que no son de polic\'ia. Ejemplo: el homicidio est\'a prohibido, pero se aplica la ley del lugar de perpetraci\'on del delito. \\
\addlinespace

\textbf{¿Qu\'e ejemplos de normas de polic\'ia dio la docente?} &
La pandemia (soluciones espec\'ificas para contratos de viaje, por ejemplo) y la crisis de 2001: corralito, pesificaci\'on de contratos en d\'olares, informatizaci\'on de salarios, \emph{default} y prohibici\'on de las quiebras. \\

\midrule
\multicolumn{2}{>{\raggedright\arraybackslash}p{\dimexpr0.92\textwidth+2\tabcolsep\relax}}{%
\textbf{Resumen:} el orden p\'ublico frena la ley elegida cuando afecta principios fundamentales del Estado, como ocurre con la compraventa de marihuana, que la Corte no autoriz\'o. Las normas de polic\'ia son m\'as r\'igidas: se aplican de forma directa e inmediata aunque haya elementos de extranjer\'ia, y no toda norma coactiva lo es. La pandemia y la crisis de 2001 ilustran su aplicaci\'on.} \\
\end{longtable}

% ------------------------------------------------------------
\chapter[Protecci\'on del d\'ebil jur\'idico]{Protecci\'on del d\'ebil jur\'idico}
% ------------------------------------------------------------

\section{El d\'ebil jur\'idico como l\'imite}

Adem\'as del orden p\'ublico y de las normas de polic\'ia, existen otros l\'imites a la voluntad, entre ellos el contrato de consumo: en el derecho argentino no se puede elegir derecho extranjero en el contrato de consumo.

Se lo presenta, en general, como un l\'imite en protecci\'on del \textbf{d\'ebil jur\'idico}: aquellos casos en que las partes no est\'an en igualdad de negociaci\'on, sea contrato de consumo, de trabajo o de seguros. En el seguro se informa cu\'al es la prima y no hay mucho para negociar; la Superintendencia aprob\'o la p\'oliza y esa es la que se aplica. La parte no tiene poder de negociaci\'on, o este es bastante limitado.

Cuando hay un d\'ebil jur\'idico aparece un l\'imite a la autonom\'ia de la voluntad, pero puede darse de dos maneras distintas. La Rep\'ublica Argentina, en materia de contratos de consumo, adopt\'o la postura de que \textbf{no se puede elegir la ley aplicable}. Se interpreta que, sea lo que sea que se elija, lo impone la otra parte: nunca se elegir\'a lo que m\'as convenga al d\'ebil, porque este eval\'ua otras cosas, otros beneficios, otras circunstancias, y no toma la protecci\'on m\'as conveniente. Hay otros sistemas jur\'idicos que, en cambio, ...

% TODO: contenido incompleto o ambiguo en las notas originales (la explicaci\'on queda interrumpida por la intervenci\'on de un alumno; la docente la retoma m\'as adelante, en el punto sobre las dos modalidades de protecci\'on)

\section{D\'ebil jur\'idico y orden p\'ublico no son lo mismo}

Ante la observaci\'on de que la normativa de defensa del consumidor se aplica de todos modos por ser de orden p\'ublico, la docente aclara que la doctrina, y tambi\'en los legisladores, suelen decir que todas las normas coactivas son de orden p\'ublico. Hay que distinguir que algo sea \textbf{obligatorio} de que sea de \textbf{orden p\'ublico}: son dos cosas distintas. Se act\'ua de una u otra manera porque algo cuestiona principios fundamentales del sistema jur\'idico, o porque se cree proteger a la parte m\'as d\'ebil.

La tutela de la parte m\'as d\'ebil es considerada por algunos de orden p\'ublico, como el derecho laboral, cuyos principios son importantes. En materia laboral algunas cuestiones se consideran fundamentales. Para el derecho laboral o el de consumo, tal principio es de orden p\'ublico, pero eso no significa que todos los contratos de consumo o todas las disposiciones de derecho laboral lo sean.

Cada especialidad considera que su materia es de orden p\'ublico o de raigambre constitucional, porque el derecho del trabajo o el del consumidor est\'an en la Constituci\'on. Pero eso no significa que la Ley de Contrato de Trabajo tenga nivel constitucional; a eso apunta el razonamiento.

Se proh\'ibe a las partes elegir la ley aplicable porque se cree proteger a la parte m\'as d\'ebil, evitando que elija una ley que le sea perjudicial. Lo mismo ocurre en materia laboral: no se deja establecer libremente la cantidad de horas de trabajo y se pone un l\'imite diario, porque se supone que la necesidad de conseguir el trabajo lleva a aceptar condiciones perjudiciales.

\begin{example}{El trabajador que prefiere dos semanas seguidas}
Quiz\'a el trabajador prefiere trabajar dos semanas seguidas sin descanso y luego tomarse los cuatro fines de semana juntos, de modo que, sumado el fin de semana, sean siete d\'ias sin trabajar. Si a \'el le conviene, ¿por qu\'e no? Cuando se dice <<no, porque es de orden p\'ublico>>, cabe preguntar qu\'e otro principio se est\'a vulnerando si la persona lo hace conscientemente y porque le conviene.

Se tiene presente la Revoluci\'on Industrial: se trabajaban 18 horas, con ni\~nos de 8 a\~nos. Por eso se limita la jornada y se otorgan vacaciones.
\end{example}

La noci\'on de orden p\'ublico y la protecci\'on del d\'ebil jur\'idico son en realidad \textbf{dos carriles distintos}. La protecci\'on del d\'ebil jur\'idico trata de que este no pacte cosas que le sean perjudiciales.

\begin{table}[H]
\centering
\caption{Orden p\'ublico y protecci\'on del d\'ebil jur\'idico: dos carriles distintos}
\small
\begin{tabularx}{\textwidth}{
    >{\raggedright\arraybackslash}p{0.20\textwidth}
    >{\raggedright\arraybackslash}X
    >{\raggedright\arraybackslash}X
}
\toprule
 & \textbf{Orden p\'ublico} & \textbf{Protecci\'on del d\'ebil jur\'idico} \\
\midrule
\textbf{Fundamento} & Principios fundamentales de la sociedad y del ordenamiento en general. & Tutelar a la parte m\'as d\'ebil: que no pacte cosas que le sean desfavorables. \\
\addlinespace
\textbf{Pregunta que se hace} & ¿La ley elegida o el acto conculca principios fundamentales del Estado? & ¿Las partes negociaron en igualdad? ¿La elecci\'on protege o perjudica al d\'ebil? \\
\addlinespace
\textbf{Ejemplos} & Objeto il\'icito; moral y buenas costumbres. & Contratos de consumo, de trabajo y de seguro. \\
\addlinespace
\textbf{Cuidado} & No todo lo obligatorio es de orden p\'ublico. & No se lo debe meter en <<la bolsa>> del orden p\'ublico. \\
\bottomrule
\end{tabularx}
\end{table}

\section{Dos modalidades de protecci\'on: prohibir o controlar la ley elegida}

El Convenio de Roma, hoy Reglamento de Roma, rige en materia de contratos civiles y comerciales y est\'a vigente en la Uni\'on Europea. Permite la autonom\'ia de la voluntad en materia de contrato de trabajo, de seguro y de consumo, siempre y cuando la ley elegida sea \textbf{m\'as favorable y m\'as protectoria} para el consumidor, el trabajador o el asegurado.

Si se pact\'o una ley y, llegado el juicio, el juez advierte que esa ley perjudica al consumidor o al trabajador, este puede pedir que no se le aplique y que se aplique su derecho propio, por ejemplo la ley del lugar de ejecuci\'on del trabajo. Si la ley pactada no es m\'as perjudicial, sino que beneficia, no hay raz\'on para dejarla de lado.

\begin{example}{El derecho del consumidor en Brasil}
En el \'ambito sudamericano, la ley m\'as protectoria en derecho de consumo es la brasile\~na. Tiene el c\'odigo de defensa de los derechos del consumidor, que exige determinado control de calidad y la extensi\'on de las responsabilidades; es mucho m\'as meticuloso y protectorio del consumidor que otras disposiciones. Si en un contrato se eligi\'o el derecho de Brasil y este da m\'as beneficios, ¿con qu\'e fundamento se dir\'ia que, por orden p\'ublico, no corresponde aplicar una ley m\'as favorable para el consumidor? Ser\'ia contradictorio.
\end{example}

% TODO: contenido incompleto o ambiguo en las notas originales (en la grabaci\'on se escucha ``una ley que es m\'as desfavorable''; por el sentido del razonamiento el apunte lo corrigi\'o a ``m\'as favorable'')

Por eso se advierte sobre meter en la ``bolsa'' del orden p\'ublico cosas que tienen otro fundamento. El fundamento del d\'ebil jur\'idico es buscar la tutela y evitar que elija cosas desfavorables; el orden p\'ublico, en cambio, habla de los principios fundamentales de la sociedad en general y del ordenamiento en general.

\subsection{Debate: ¿cu\'al modalidad protege m\'as?}

Los l\'imites de protecci\'on a los d\'ebiles jur\'idicos operan, entonces, de dos formas. Una es \textbf{prohibir directamente} la elecci\'on, que es lo que hace la Argentina. La otra es \textbf{permitir la autonom\'ia de la voluntad} en tanto y en cuanto el derecho elegido sea m\'as favorable a los intereses del d\'ebil jur\'idico: se permite, pero hay un control de calidad respecto de la ley elegida, y si perjudica, no se aplica. Se plante\'o que esta segunda opci\'on lleva a una protecci\'on mayor.

La docente observ\'o que se est\'a acostumbrado a que no se puede, sin preguntarse por qu\'e no se puede elegir algo mejor. Una cosa es decir que no se puede elegir porque se protege al d\'ebil jur\'idico, y otra, que se pueda elegir en tanto se vea que la ley lo favorece. Hay normas que van a la prohibici\'on por la prohibici\'on misma.

Frente a la objeci\'on de que quien contrata por necesidad (la compra de algo necesario, un trabajo aceptado ``por el pan en la mesa'') acepta aunque no le convenga, se respondi\'o que se dicta entonces una ley que lo proteja y que no se puede tocar, con lo que se ahorra tener que evaluar. La docente agreg\'o que es m\'as f\'acil presumirlo que evaluarlo caso por caso, que tomarse el trabajo de dejar optar. Y que despu\'es, en todo caso, el trabajador puede reclamar que se le aplique la norma que le corresponde, porque la que se eligi\'o no cumple lo que se hab\'ia pactado (por ejemplo, haber pactado trabajar sin feriados para tomarlos todos juntos y que luego no se los otorguen as\'i).

% TODO: contenido incompleto o ambiguo en las notas originales (este tramo de la grabaci\'on tiene superposici\'on de voces; la distribuci\'on de intervenciones es la m\'as coherente con el sentido del di\'alogo)

Los l\'imites dependen de la extensi\'on que se quiera dar a la autonom\'ia de la voluntad. Si se permite elegir en tanto la ley elegida beneficie, la autonom\'ia es mucho m\'as extensa. Decir <<no se puede elegir en contrato de consumo, de trabajo y de seguro; no se puede, punto>> es un l\'imite liso y llano: ah\'i no hay autonom\'ia. En cambio, si el Estado permite elegir siempre que la ley elegida beneficie, la autonom\'ia de las partes es m\'as amplia, sujeta despu\'es a una evaluaci\'on de control de calidad de las leyes elegidas.

Se tienen as\'i tres l\'imites: orden p\'ublico, norma de polic\'ia y protecci\'on del d\'ebil jur\'idico.

\begin{table}[H]
\centering
\caption{Dos modalidades de protecci\'on del d\'ebil jur\'idico}
\small
\begin{tabularx}{\textwidth}{
    >{\raggedright\arraybackslash}p{0.20\textwidth}
    >{\raggedright\arraybackslash}X
    >{\raggedright\arraybackslash}X
}
\toprule
 & \textbf{Prohibici\'on de elegir la ley} & \textbf{Elecci\'on permitida si la ley es m\'as favorable} \\
\midrule
\textbf{C\'omo funciona} & No se puede elegir la ley aplicable. <<Punto.>> & Se puede elegir, pero hay un <<control de calidad>> de la ley elegida. \\
\addlinespace
\textbf{Qui\'en la adopta (seg\'un la clase)} & Rep\'ublica Argentina, en contratos de consumo. & Reglamento de Roma (UE): trabajo, seguro y consumo. \\
\addlinespace
\textbf{Extensi\'on de la autonom\'ia} & M\'inima: <<l\'imite liso y llano>>. & M\'as amplia, sujeta a evaluaci\'on de la ley elegida. \\
\addlinespace
\textbf{Ventaja se\~nalada} & Es m\'as f\'acil: se presume el perjuicio y se <<ahorra>> evaluar caso por caso. & Si la ley elegida beneficia al d\'ebil, no hay raz\'on para sac\'arsela. \\
\bottomrule
\end{tabularx}
\end{table}

\section{Cuadro Cornell}

\begin{longtable}{
    >{\raggedright\arraybackslash}p{0.26\textwidth}
    >{\raggedright\arraybackslash}p{0.66\textwidth}
}
\toprule
\textbf{Preguntas / palabras clave} & \textbf{Notas / respuestas} \\
\midrule
\endfirsthead
\toprule
\textbf{Preguntas / palabras clave} & \textbf{Notas / respuestas} \\
\midrule
\endhead
\bottomrule
\endfoot
\bottomrule
\endlastfoot

\textbf{¿Qui\'en es el d\'ebil jur\'idico y por qu\'e limita la autonom\'ia?} &
La parte sin igualdad de negociaci\'on: consumidor, trabajador, asegurado. Se limita la autonom\'ia para que no pacte cosas que le sean perjudiciales. \\
\addlinespace

\textbf{¿Orden p\'ublico y d\'ebil jur\'idico son lo mismo?} &
No: son <<dos carriles distintos>>. Algo obligatorio no es necesariamente de orden p\'ublico. El d\'ebil jur\'idico busca tutela; el orden p\'ublico protege los principios fundamentales del ordenamiento en general. \\
\addlinespace

\textbf{¿C\'omo se protege al d\'ebil jur\'idico?} &
De dos maneras: prohibiendo elegir la ley (Argentina, consumo) o permitiendo elegirla si es m\'as favorable (Reglamento de Roma). La segunda es m\'as amplia, sujeta a un control de calidad de la ley elegida. \\

\midrule
\multicolumn{2}{>{\raggedright\arraybackslash}p{\dimexpr0.92\textwidth+2\tabcolsep\relax}}{%
\textbf{Resumen:} la protecci\'on del d\'ebil jur\'idico se distingue del orden p\'ublico y puede operar por prohibici\'on de elegir la ley o por control de la ley m\'as favorable. La primera ofrece una soluci\'on m\'as simple; la segunda concede una autonom\'ia m\'as amplia, sujeta a evaluaci\'on. Los l\'imites dependen de la extensi\'on que cada Estado quiera dar a la autonom\'ia de la voluntad.} \\
\end{longtable}

% ------------------------------------------------------------
\chapter[Otros l\'imites y momento de la elecci\'on]{Otros l\'imites y momento de la elecci\'on de ley}
% ------------------------------------------------------------

\section{Moral, buenas costumbres y elementos de extranjer\'ia}

El primero y principal de estos l\'imites es la \textbf{moral y las buenas costumbres}, que entran dentro de los principios del orden p\'ublico, en la ``bolsa general'' del orden p\'ublico.

Otro l\'imite fundamental, y el principal en cuanto a la posibilidad de elegir ley aplicable, es que el contrato tenga \textbf{elementos de extranjer\'ia}: elementos que vinculen distintos sistemas jur\'idicos.

\begin{example}{El contrato sin ning\'un elemento extranjero}
El contrato se celebra entre dos argentinos, comprador y vendedor, en la ciudad de Buenos Aires; el lugar de cumplimiento, la entrega de la mercader\'ia, es en Corrientes, Rep\'ublica Argentina; el garante es la madre del comprador, que es argentina; se contrata una compa\~n\'ia de seguro argentina. ¿Para qu\'e aplicar el derecho de Rusia? Si el contrato tiene todos los elementos reunidos o conectados con un \'unico ordenamiento jur\'idico, no se puede elegir derecho extranjero: se aplica la normativa nacional de ese Estado.
\end{example}

Para poder hacer autonom\'ia conflictual, el contrato debe tener estos elementos de extranjer\'ia: debe ser internacional o, en el caso de la clase, regional, es decir, vincular de alguna manera a los pa\'ises del Mercosur. Por ejemplo, el domicilio del comprador est\'a en un Estado y el del vendedor en otro, aunque el contrato se haya celebrado y ejecutado en el mismo: hay un elemento de extranjer\'ia.

\begin{important}
Sin elementos de extranjer\'ia no hay autonom\'ia conflictual: si todos los elementos del contrato est\'an conectados con un \'unico ordenamiento, se aplica la normativa nacional de ese Estado. Basta con un elemento de extranjer\'ia, por ejemplo el domicilio de cada parte en un Estado distinto.
\end{important}

\section{Momento de la elecci\'on: autonom\'ia amplia o cerrada}

Otro l\'imite es el tiempo en el cual se puede elegir la ley aplicable. Una \textbf{autonom\'ia amplia} permitir\'ia elegir la ley aplicable al contrato en cualquier momento de la relaci\'on jur\'idica, porque el contrato es el acuerdo de voluntades de dos partes. Si ambas est\'an de acuerdo en cambiar la ley aplicable y no perjudican a nadie ---por ejemplo, aplicar la ley argentina en lugar de la de Uruguay---, lo deciden ellas y no molestan a nadie; no hay cuestiones de orden p\'ublico en juego.

Una \textbf{autonom\'ia m\'as cerrada} permitir\'ia hacer uso de la autonom\'ia de la voluntad solo al momento de celebrar el contrato. Una vez celebrado, queda establecida la ley aplicable, ya sea porque las partes eligieron o porque no eligieron; si no eligieron, lo determina la norma de conflicto. Pero la ley aplicable queda fijada y no se puede modificar despu\'es.

\begin{table}[H]
\centering
\caption{Autonom\'ia amplia y autonom\'ia cerrada en el tiempo}
\small
\begin{tabularx}{\textwidth}{
    >{\raggedright\arraybackslash}p{0.22\textwidth}
    >{\raggedright\arraybackslash}X
    >{\raggedright\arraybackslash}X
}
\toprule
 & \textbf{Autonom\'ia amplia} & \textbf{Autonom\'ia cerrada} \\
\midrule
\textbf{Cu\'ando se puede elegir o cambiar la ley} & En cualquier momento de la relaci\'on jur\'idica. & Solo al celebrar el contrato. \\
\addlinespace
\textbf{Si no se eligi\'o} & Pueden elegirla o cambiarla despu\'es. & Rige lo que diga la norma de conflicto; queda fijada y no se modifica. \\
\addlinespace
\textbf{Raz\'on} & El contrato es acuerdo de voluntades: si ambas partes est\'an de acuerdo y no perjudican a nadie, no hay problema. & Seguridad: la ley aplicable queda determinada desde la celebraci\'on. \\
\bottomrule
\end{tabularx}
\end{table}

\subsection{Condiciones l\'ogicas para una autonom\'ia amplia}

Para que una autonom\'ia amplia sea l\'ogica, se se\~nalaron condiciones. La primera es que \textbf{no afecte derechos de terceros}: la nueva ley no puede afectarlos. Por ejemplo, si el garante ten\'ia una determinada responsabilidad limitada a cierta extensi\'on y, al cambiar la ley aplicable, pasa a responder m\'as, se afecta un derecho por una relaci\'on pactada entre comprador y vendedor. El derecho elegido no puede afectar los intereses de terceros ni la validez del contrato: si seg\'un la nueva ley el contrato resultara nulo, no tendr\'ia sentido. No pueden alterarse ni la validez formal del contrato ni los derechos adquiridos por terceros.

Respecto de hasta cu\'ando puede hacerse el cambio, se respondi\'o que \textbf{hasta el cumplimiento}: una vez cumplido el contrato, ya no se puede. Ante la pregunta de si podr\'ia elegirse la ley aplicable durante el juicio, un alumno respondi\'o que durante el juicio no, pero que antes de ir a juicio s\'i se puede pactar, como otra autonom\'ia de la voluntad.

\begin{important}
La nueva ley elegida no puede afectar (1) los derechos adquiridos por terceros ni (2) la validez formal del contrato. En el debate, los alumnos ubicaron el l\'imite temporal en el cumplimiento del contrato: una vez cumplido, ya no se puede cambiar la ley.
\end{important}

\section{El Reglamento de Roma y la elecci\'on t\'acita de la ley}

% TODO: contenido incompleto o ambiguo en las notas originales (en este tramo final la grabaci\'on no permite identificar con certeza a cada interviniente y termina durante la discusi\'on; la atribuci\'on de los turnos es la m\'as coherente con el sentido del di\'alogo)

El Reglamento de Roma sobre obligaciones contractuales permite la autonom\'ia y el cambio de ley aplicable \textbf{en cualquier momento, aun durante el proceso}, y tanto de forma escrita como de forma t\'acita. Es decir, se puede modificar, durante el proceso y de com\'un acuerdo con la otra parte, la ley que el juez deb\'ia aplicar, y hacerlo de manera t\'acita.

Se plante\'o que esto generar\'ia inseguridad (``a nosotros nos dar\'ia p\'anico''), y que corresponde al juez evitar que se viole la seguridad jur\'idica entre las partes; tambi\'en se pregunt\'o c\'omo se sabe que hubo acuerdo t\'acito. La docente respondi\'o que hay elecci\'on t\'acita cuando una parte empieza a invocar derecho extranjero de un tercer Estado y la otra parte contesta sobre la misma normativa: eso es lo que Roma define como elecci\'on t\'acita.

\begin{definition}{Elecci\'on t\'acita de la ley (Reglamento de Roma)}
Seg\'un lo expuesto en clase, la ley aplicable puede elegirse o cambiarse en cualquier momento ---antes, durante o despu\'es de surgido el conflicto, incluso durante el juicio--- y de forma expresa o t\'acita. Hay elecci\'on t\'acita cuando una parte invoca el derecho de un Estado y la otra contesta sobre la misma normativa.
\end{definition}

\subsection{Objeci\'on planteada}

Un alumno objet\'o que, si la parte invoca por escrito en un expediente judicial determinada norma de un c\'odigo extranjero, firmada por ella y por su abogado, y la otra contesta por escrito invocando otra norma del mismo c\'odigo, ``lo estamos poniendo por escrito: ¿qu\'e t\'acito?''. La docente respondi\'o que ese es precisamente el ejemplo que da Roma cuando dice que se puede hacer en cualquier momento, antes, durante o despu\'es de surgido el conflicto, aun durante el juicio, y de forma expresa o t\'acita: es decir, qu\'e amplia es la autonom\'ia de la voluntad.

El alumno insisti\'o en que, cuando se define lo t\'acito como el caso en que un derecho es invocado por una parte y la otra contesta, se interpreta que la cl\'ausula escrita de ley aplicable pactada en el contrato fue cambiada de esa forma. A su juicio no ser\'ia una forma t\'acita, sino un \textbf{acuerdo no simult\'aneo}: uno lo hizo en un momento y el otro confirm\'o la elecci\'on de ese derecho en otro; es expreso y no t\'acito. Ser\'ia m\'as t\'acito, por ejemplo, enviar una orden de compra sin firmar.

\section{Cuadro Cornell}

\begin{longtable}{
    >{\raggedright\arraybackslash}p{0.26\textwidth}
    >{\raggedright\arraybackslash}p{0.66\textwidth}
}
\toprule
\textbf{Preguntas / palabras clave} & \textbf{Notas / respuestas} \\
\midrule
\endfirsthead
\toprule
\textbf{Preguntas / palabras clave} & \textbf{Notas / respuestas} \\
\midrule
\endhead
\bottomrule
\endfoot
\bottomrule
\endlastfoot

\textbf{¿Qu\'e otros l\'imites mencion\'o la docente?} &
La moral y las buenas costumbres (dentro del orden p\'ublico) y la necesidad de elementos de extranjer\'ia: sin ellos no se puede elegir derecho extranjero. \\
\addlinespace

\textbf{¿Cu\'ando se puede elegir o cambiar la ley aplicable?} &
Autonom\'ia amplia: en cualquier momento. Autonom\'ia cerrada: solo al celebrar el contrato. Si se cambia, no puede afectar derechos de terceros ni la validez formal del contrato; los alumnos se\~nalaron como l\'imite el cumplimiento. \\
\addlinespace

\textbf{¿Qu\'e dice el Reglamento de Roma sobre el momento y la forma de la elecci\'on?} &
Permite elegir o cambiar la ley en cualquier momento, incluso durante el proceso, de forma expresa o t\'acita. Hay elecci\'on t\'acita cuando una parte invoca el derecho de un Estado y la otra contesta sobre la misma normativa. \\
\addlinespace

\textbf{¿Qu\'e objeci\'on se plante\'o a la elecci\'on t\'acita?} &
Que, si lo invocado consta por escrito en el expediente, el acuerdo es expreso, aunque no simult\'aneo. M\'as t\'acito ser\'ia, por ejemplo, mandar una orden de compra sin firmar. \\

\midrule
\multicolumn{2}{>{\raggedright\arraybackslash}p{\dimexpr0.92\textwidth+2\tabcolsep\relax}}{%
\textbf{Resumen:} para elegir ley el contrato debe tener elementos de extranjer\'ia, y la moral y las buenas costumbres integran el orden p\'ublico. Respecto del momento, la autonom\'ia puede ser amplia (en cualquier momento) o cerrada (solo al celebrar), con l\'imites frente a terceros y la validez del contrato. El Reglamento de Roma admite la elecci\'on en cualquier momento, incluso durante el proceso, expresa o t\'acita. La libertad de las partes es real pero nunca absoluta: cada Estado decide cu\'anta le concede y d\'onde pone el freno.} \\
\end{longtable}

\end{document}
